Report the findings in the order needed to answer the research question. Distinguish measured or derived results from interpretation, and give uncertainty estimates and sample sizes where appropriate.

\begin{table}[htbp]
\centering
\caption{Example layout for a compact results table. Replace all instructional cells.}
\label{tab:results-summary}
\begin{tabularx}{\linewidth}{@{}P{0.28\linewidth}YY@{}}
\toprule
\paperTableHead{Outcome or comparison} & \paperTableHead{Result} & \paperTableHead{Supporting evidence} \\
\midrule
Primary outcome & Insert the verified result. & Specify the analysis and uncertainty. \\
Secondary outcome & Insert the verified result. & Specify the analysis and uncertainty. \\
Validation result & Insert the verified result. & Specify the control or validation setting. \\
\bottomrule
\end{tabularx}
\end{table}

Table~\ref{tab:results-summary} demonstrates the compact table style. Replace its instructional text with results from the study and keep each table tied to a specific finding.

\begin{table}[htbp]
\centering
\caption{Numeric alignment example. Values below demonstrate typesetting only; replace them with verified study results or remove this table.}
\label{tab:numeric-example}
\begin{tabularx}{\linewidth}{@{}Y S[table-format=3.0] S[table-format=1.3(3)] S[table-format=1.3]@{}}
\toprule
\paperTableHead{Comparison} & {\paperTableHead{Count}} & {\paperTableHead{Estimate $\pm$ SD}} & {\paperTableHead{Score}} \\
\midrule
Example A & 120 & 0.385 +- 0.006 & 0.759 \\
Example B & 98 & 0.608 +- 0.051 & 0.206 \\
\bottomrule
\end{tabularx}
\paperTableNote{Notes. Define units, uncertainty, missing values, and the meaning of emphasis here. Keep long methodological explanations out of column headings.}
\end{table}

\begin{figure}[htbp]
\centering
\begin{subfigure}[t]{0.48\linewidth}
  \centering
  \paperFigurePlaceholder{Insert a vector figure\\with readable axis labels.}
  \caption{First comparison.}
  \label{fig:example-a}
\end{subfigure}\hfill
\begin{subfigure}[t]{0.48\linewidth}
  \centering
  \paperFigurePlaceholder{Insert a companion figure\\at the same visual scale.}
  \caption{Second comparison.}
  \label{fig:example-b}
\end{subfigure}
\caption{\textbf{Example of a two-panel figure.} Replace the panels with the study's own graphics. State what each panel shows and explain symbols, units, sample sizes, and uncertainty in the caption.}
\label{fig:example-panels}
\end{figure}

\subsection{Primary findings}
Describe the findings that address the main objective. Connect the text to the relevant figures or tables and report the evidence needed to assess the result.

\subsection{Additional analyses}
Present relevant secondary analyses, comparisons, or validation results. Explain their relationship to the primary findings and avoid repeating the contents of tables.
